%% ---------------------------------------------------------------------
\section{Introduction and hardware in context}
\label{sec:intro}

\subsection{The GPS-denied navigation problem}

Small multirotor UAVs estimate their state by fusing a low-cost micro-electro-mechanical
(MEMS) inertial measurement unit (IMU) with GNSS position and velocity
\cite{groves2013principles}. GNSS is routinely unavailable or untrustworthy: signals
are blocked in mines, tunnels and buildings, degraded in urban canyons and beneath
foliage, and can be deliberately jammed or spoofed \cite{humphreys2008spoofing}.
When the absolute fix disappears, the vehicle must dead-reckon on the IMU, and the
error of a strapdown solution grows rapidly \cite{titterton2004strapdown}. For an
uncompensated accelerometer bias $\mathbf{b}_a$, double integration gives a position
error
\begin{equation}
  \delta\mathbf{p}(t) \;\approx\; \tfrac{1}{2}\,\mathbf{b}_a\,t^{2},
  \label{eq:bias-drift}
\end{equation}
while a small attitude error $\delta\theta$ projects gravity into the horizontal
channel, adding $\tfrac12 g\,\delta\theta\,t^2$ \cite{titterton2004strapdown}. A residual bias of
only \SI{1}{\milli\gacc} (\SI{9.8e-3}{\metre\per\second\squared}) produces roughly
\SI{0.5}{\metre} of error after \SI{10}{\second} and \SI{18}{\metre} after one minute. \Cref{sec:val-dr} confirms this growth on data
from the prototype's own autopilot. A GPS-denied navigation system must therefore
provide an \emph{exteroceptive} reference (vision, range, radio or field-map matching)
that bounds this drift, together with a means of detecting and avoiding obstacles in
unknown surroundings.

\subsection{Formal statement}

Following our problem statement, the vehicle state at time $t$ is
\begin{equation}
  \mathbf{x}_t = \left[\,\mathbf{p}_t,\;\mathbf{v}_t,\;\mathbf{q}_t,\;\mathbf{b}_a,\;\mathbf{b}_g\,\right],
\end{equation}
comprising position, velocity, attitude quaternion, and accelerometer and gyroscope
biases. Given control/IMU inputs $\mathbf{u}_t$ and exteroceptive measurements
$\mathbf{z}_{t,k}$ (camera features, LiDAR ranges, flow, rangefinder), the
maximum-a-posteriori trajectory is
\begin{equation}
  \mathbf{x}^{*}_{0:T} = \arg\max_{\mathbf{x}_{0:T}}\;
  p(\mathbf{x}_0)\prod_{t=1}^{T} p(\mathbf{x}_t\mid\mathbf{x}_{t-1},\mathbf{u}_t)
  \prod_{k} p(\mathbf{z}_{t,k}\mid\mathbf{x}_t),
  \label{eq:map}
\end{equation}
which is solved recursively by an extended Kalman filter (EKF) or in batch by
factor-graph optimisation \cite{dellaert2017factor}. In the present prototype the
recursive solution is delegated to the autopilot's EKF3 where it is available, and a
lightweight dead-reckoning front-end on the ground computer supplies the planner when
it is not (\cref{sec:design-odom}).

\subsection{Target performance}

We set the key performance indicators (KPIs) in \cref{tab:kpi}. They are
\emph{targets for the final onboard system}. \Cref{sec:discussion} reports how far the
current prototype has been verified against each.

\begin{table}[H]
\centering\small
\caption{Project KPIs (from \code{docs/PROBLEM_STATEMENT.md}).}
\label{tab:kpi}
\begin{tabularx}{\textwidth}{l l L}
\toprule
\textbf{Metric} & \textbf{Target} & \textbf{Typical fallback without GNSS} \\
\midrule
Trajectory error & $<0.5\%$ of distance flown & $>10\%$ \\
Hover drift & $\le\pm\SI{5}{\centi\metre}$ over \SI{10}{\minute} & $\ge\SI{2}{\metre}$ per minute \\
Update rate & state $\ge\SI{100}{\hertz}$, vision $\ge\SI{30}{\hertz}$ & dropped frames break the filter \\
Compute power & $\le\SI{15}{\watt}$, onboard & needs off-board computing \\
Illumination & \SIrange{0}{10000}{lux} and above & fails below \SI{50}{lux} \\
Endurance & $\ge\SI{30}{\minute}$ (objective 8) & -- \\
\bottomrule
\end{tabularx}
\end{table}

\subsection{Background and related work}

\paragraph{Exteroceptive references} Visual and visual--inertial odometry and SLAM
(e.g.\ VINS-Mono \cite{qin2018vins}, ORB-SLAM3 \cite{campos2021orbslam3}) achieve
sub-percent drift with light, inexpensive cameras, but fail in darkness, smoke and on
textureless surfaces. Downward optical-flow sensors with a rangefinder
\cite{honegger2013px4flow} give metric velocity at negligible compute cost but no map.
2D and 3D LiDAR provide illumination-independent geometry and support occupancy-grid
SLAM \cite{macenski2021slamtoolbox} at a weight and power cost. Radio beacons (UWB),
fiducials, magnetic-anomaly (MAGNAV) and terrain-referenced (TERCOM) navigation need
prior infrastructure or maps. Surveys of the space include \cite{balamurugan2016survey}.

\paragraph{Open robotics middleware} ROS~2 \cite{macenski2022ros2} with the Nav2
navigation stack \cite{macenski2020marathon} offers costmaps, global planners, local
controllers such as the dynamic window approach \cite{fox1997dwa}, and behaviour-tree
mission logic \cite{colledanchise2018bt}. PX4 \cite{meier2015px4} and ArduPilot
\cite{ardupilot} are the dominant open autopilots, communicating over MAVLink
\cite{koubaa2019mavlink}.

\paragraph{Gap} Integrated, documented, reproducible GPS-denied systems for hobby-grade
multirotors are rare. Commercial modules with onboard VIO exist but are closed, costly
or discontinued. Academic systems are typically tied to bespoke airframes and high-end
compute. There is a practical need for an \emph{incremental} path: a system that can be
assembled from parts already found in student laboratories, that exposes every signal
for inspection, and that can migrate towards fully onboard processing.

\subsection{Contribution}

Our main contributions are:
\begin{enumerate}[leftmargin=*,itemsep=1pt]
  \item a structured trade study of 15 GPS-denied navigation modalities with
    multi-criteria scoring for three mission profiles (\cref{sec:design-trade});
  \item a companion-computer architecture that adds 2D LiDAR, optical-flow odometry,
    Nav2 planning and semantic perception to an unmodified ArduPilot vehicle, using
    only USB connections on the airframe (\cref{sec:arch});
  \item a latency-engineered Wi-Fi link with single-datagram encodings, QoS marking,
    clock synchronisation and per-stream health metrics (\cref{sec:design-link});
  \item a layered command-safety design: arming and mode gating, speed clamping,
    stale-command stop, a vibration guard and pilot override (\cref{sec:design-safety});
  \item a hybrid YOLO--SAM segmentation module with explicit mask-quality gates and
    unit tests (\cref{sec:design-perception});
  \item a reproducible PX4/Gazebo simulation environment, including a generated office
    world and frontier exploration (\cref{sec:op-sim});
  \item a clear account of what we have and have not yet verified
    (\cref{sec:validation,sec:discussion}).
\end{enumerate}

\subsection{Prototype versus design target}
\label{sec:proto-vs-target}

We describe two hardware configurations, and it is worth keeping them apart:
\begin{description}[leftmargin=1em,itemsep=2pt]
  \item[Current prototype (implemented).] F450 quadrotor with a Pixhawk~6C running ArduPilot, YDLIDAR G2, USB
    camera and autopilot-attached optical flow/rangefinder. A Raspberry Pi streams the
    sensor data to a GPU laptop, which performs estimation, planning and perception.
    All software in \cref{sec:arch,sec:design,sec:build,sec:operation} refers to this
    configuration.
  \item[Design target (specified, not yet built).] A SC132GS stereo camera,
    Raspberry Pi~5 (8\,GB), Holybro Kakute H7 running PX4 and an ELRS radio link, with
    \emph{all} processing onboard (\cref{tab:bom-target}). Our timeline
    (\cref{tab:timeline}) places this in phases~1--6.
\end{description}
Off-board processing is not part of the final design. We use the ground laptop as a
development step, because it lets us watch every intermediate signal while the
algorithms are still being tuned.
